\section{Related Work}
\label{sec:related}

\subsection{Backbones for Irregularly Sampled Clinical Time Series}

The modelling substrate for early warning is a sequence of observations whose sampling times carry as much information as their values. The most recent work in this line treats irregularity as a first-class object rather than an inconvenience: continuous-time formulations parameterise the latent trajectory between observations and evaluate it at arbitrary query times \citep{shen2026enhancing}, reliability-aware aggregation weights observations by their estimated trustworthiness before temporal pooling \citep{kim2026retamamba}, and graph constructions connect observations across variables and time so that forecasting can propagate information along non-uniform edges \citep{yalavarthi2024grafiti}. A complementary strand models the observation process itself, on the argument that the decision to order a laboratory test is a clinical action and therefore informative; this idea has been used both to synthesise realistic irregular trajectories \citep{mehdizavareh2026informative} and to build masked-attention predictors for deterioration on general wards \citep{mun2026made}. Domain-specific instantiations extend the same machinery to data-scarce settings such as venous thromboembolism detection after cancer treatment \citep{xu2026bridging}, and surveys of the area converge on the observation that architectural progress has outpaced agreement on evaluation \citep{sun2026review}. We adopt these backbones without modification and treat the score function they produce as an input to our policy layer, which is where our contribution lies.

\subsection{Early Warning and Deterioration Prediction}

Clinical early warning has moved from static severity scores toward dynamic, trajectory-aware prediction. Recent multicentre work predicts the shape of sepsis deterioration rather than a single risk estimate, arguing that heterogeneous trajectories require patient-specific forecasting \citep{xiumin2026machine}, while efficiency-oriented designs pursue compact models that remain explainable under the compute constraints of a ward workstation \citep{taparugssanagorn2026lightweight}. Others inject clinical reasoning into the prediction pathway: dynamic temporal graphs embed calculator-derived features so that the model inherits the semantics of established instruments \citep{yin2024sepsiscalc}, and language-model-guided simulation is used to generate interpretable physiological dynamics for early warning \citep{nie2026clinically}. Transfer-oriented formulations address the practical reality that a model trained in one institution must be reused in another with different coding practice and case mix \citep{komolafe2025early}, and earlier work established that deep models over longitudinal inpatient records can outperform tabular baselines for outcome prediction \citep{rong2025deep,yong2024deep,mohammad2024early}. Across this literature the model is the unit of contribution; the alerting rule that converts its output into an operational signal is treated as an afterthought, and comparison between methods is almost always conducted at a fixed threshold rather than at a fixed alarm volume.

\subsection{Calibration, Selective Prediction, and Conformal Guarantees}

A parallel literature asks not whether a model is accurate but whether its outputs can be trusted at a stated level. Systematic analysis of calibration and selective prediction for multimodal clinical classification shows that abstention behaviour, not raw discrimination, determines whether a system is safe to deploy \citep{lpez2026empirical}. For time-series classifiers, reliability estimation has been extended beyond post hoc remapping of output scores toward evidence that a confident prediction is actually supported by the input \citep{cenacchi2026beyond}, and audits of adaptation procedures have exposed failure modes that only appear when the update rule is evaluated with realistic label delays \citep{fentazi2026training}. Conformal prediction supplies the distribution-free coverage guarantee that these approaches seek, and has been applied to physiological regression where calibrated uncertainty is a safety requirement \citep{strickland2025degrees}. Adaptation at inference time has also been used to counter patient-specific shift in haemodynamic prediction \citep{li2026crosssample}. What this body of work assumes, and what an operational alarm system violates, is independence: a decision rule with a refractory period produces a dependent sequence, so guarantees derived from exchangeable examples do not transfer to the alarm rate actually experienced by a clinical unit. Our procedure is designed for exactly this dependent regime.

\subsection{Evaluation Practice}

Finally, the reporting conventions of the field shape what gets optimised. Risk-of-bias instruments for prediction models have recently been extended to cover artificial-intelligence-specific failure modes, including the handling of calibration, heterogeneity, and operational applicability \citep{moons2025probastai}, and methodological reviews of intensive-care adverse-event prediction note that alert burden and time-to-alert remain inconsistently reported \citep{wang2026deep,hernandezdiaz2026deep,swinckels2024deep}. We take these critiques as a specification: every experiment below reports an operational quantity alongside a statistical one, and all comparisons between alerting strategies are conducted at matched alarm budgets so that discrimination is never bought with unreported alarm volume.
